Most people are familiar with grammar in the context of natural language.
In English class we learn that we shouldn't end a sentence with a preposition.
This is an example of a \textit{grammatical rule}.
In broad strokes, grammar refers to the rules that govern what constitutes a well-formed sentence in a languange.
It's the concern of grammar to determine that \enquote{I taught the car in the refrigerator} is well-formed, whereas \enquote{I car teach refrigerator in there} is not.
It's \textit{not} the concern of grammar to determine that the first sentence has no reasonable interpretation in English.

The grammar of a programming language determines what constitutes a well-formed \textit{program} in that language.
Due to the grammatical precision of programming languages, these tend to be called \textbf{formal grammars}.
In OCaml, the program
%
\begin{ocaml}
  let f x = x + 1
\end{ocaml}
%
is well-formed, but the program
%
\begin{ocaml}
  f let x = x 1 +
\end{ocaml}
%
is not.
As with natural language, the grammar of a programming language is not concerned with the \textit{meaning} of programs, just their well-formedness, e.g., the program
%
\begin{ocaml}
  let omega x = x x
\end{ocaml}
%
is well-formed, but isn't well-typed.

A good programming language has a clear specification of its grammar.
The grammar of OCaml, for example, is given in its entirety in \href{https://ocaml.org/manual/5.3/expr.html}{The OCaml Manual}.
After going though this chapter, you should be able to read the specification given there.
\bigskip

\todo{Goals of the chapter.}

\section{A Thought Experiment}

How do we know that this sentence is grammatically correct?
Let's try to break down the cognitive process of this determination.
Consider the following English statement.
\begin{align*}
  \newcommand{\tm}{\bnft}
  \tm{the} \ \tm{cow} \ \tm{jumped} \ \tm{over} \ \tm{the} \ \tm{moon}
\end{align*}
We might first recognize that each word is a particular part of speech (noun, verb, preposition, etc.) and that, with regards to grammar, only the part of speech counts.
We represent this step of the process by replacing each word with an abstract symbol that stands in for its part of speech.
\begin{align*}
  \newcommand{\nt}{\bnfnt}
  \nt{article} \ \nt{noun} \ \nt{verb} \ \nt{prep} \ \nt{article} \ \nt{noun}
\end{align*}
After this, there are some familiar patterns: $\bnfnt{article} \ \bnfnt{noun}$ is the determination or quantification of an object, e.g., \enquote{the house} or \enquote{a car.}
In other words, this pattern represents a \textit{noun phrase} or \textit{nominal phrase}.
We mentally group these forms, and this group can be represented by a new symbol.
\begin{align*}
  \bnfnt{noun-phrase} \ \bnfnt{verb} \ \bnfnt{prep} \ \bnfnt{noun-phrase}
\end{align*}
Another pattern we might recognize: a proposition followed by a noun phrase is a single unit, e.g., \enquote{over the moon}, \enquote{through the woods}, and \enquote{behind the wall.}
These are called \textit{prepositional phrases} and can, again, be represented by a new symbol.
\begin{align*}
  \newcommand{\nt}{\bnfnt}
  \nt{noun-phrase} \ \nt{verb} \ \nt{prep-phrase}
\end{align*}
Prepositional phrases modify verbs, creating a \textit{verb phrase} (e.g., \enquote{ran to the car}) leaving us with something like:
\begin{align*}
  \newcommand{\nt}{\bnfnt}
  \nt{noun-phrase} \ \nt{verb-phrase}
\end{align*}
which we should finally recognized as the canonical structure of a well-formed sentence: \textit{thing does}.

There are general structures that are allowed by a grammar, and these structures can be \textit{instantiated} at particular words in the language.
The rules that define these structures are often heirarchical; a sentence is made of parts (e.g, a noun phrase and a verb phrase) and those constituent parts are made of smaller parts (e.g., an article and a noun) and so on until we reach units of our language, the particular words.
Demonstrating that a sentence is grammatically correct means demonstrating that it fits into this heirarchy of allowed structures.
If we reverse the order of steps we took above, and break up some of the steps we did in parallel, we get a \textit{proof} that our sentence is grammatically correct.
%
\begin{align*}
  & \bnfnt{sentence} \\
  & \bnfnt{noun-phrase} \ \bnfnt{verb-phrase} \\
  & \bnfnt{noun-phrase} \ \bnfnt{verb} \ \bnfnt{prep-phrase} \\
  & \bnfnt{noun-phrase} \ \bnfnt{verb} \ \bnfnt{prep} \ \bnfnt{noun-phrase} \\
  & \bnfnt{noun-phrase} \ \bnfnt{verb} \ \bnfnt{prep} \ \bnfnt{article} \ \bnfnt{noun} \\
  & \bnfnt{article} \ \bnfnt{noun} \ \bnfnt{verb} \ \bnfnt{prep} \ \bnfnt{article} \ \bnfnt{noun} \\
  & \bnft{the} \ \bnfnt{noun} \ \bnfnt{verb} \ \bnfnt{prep} \ \bnfnt{article} \ \bnfnt{noun} \\
  & \bnft{the} \ \bnft{cow} \ \bnfnt{verb} \ \bnfnt{prep} \ \bnfnt{article} \ \bnfnt{noun} \\
  & \bnft{the} \ \bnft{cow} \ \bnft{jumped} \ \bnfnt{prep} \ \bnfnt{article} \ \bnfnt{noun} \\
  & \bnft{the} \ \bnft{cow} \ \bnft{jumped} \ \bnft{over} \ \bnfnt{article} \ \bnfnt{noun} \\
  & \bnft{the} \ \bnft{cow} \ \bnft{jumped} \ \bnft{over} \ \bnft{the} \ \bnfnt{noun} \\
  & \bnft{the} \ \bnft{cow} \ \bnft{jumped} \ \bnft{over} \ \bnft{the} \ \bnft{moon}
\end{align*}
%
In what follows we call this sort of proof a \textbf{derivation}.
Each line corresponds to the application of a grammatical rule.
For example, the fourth line follows from the third by way of applying the grammatical rule that a prepositional phrase can be made up of a preposition followed by a verb phrase.
If we squint, we can also start to see the tree-like structure encoded in this derivation (which we will call a \textbf{parse tree}).
\bigskip

\Tree
  [.$\bnfnt{sentence}$
    [.$\bnfnt{noun-phrase}$
      [.$\bnfnt{article}$
        $\bnft{the}$
      ]
      [.$\bnfnt{noun}$
        $\bnft{cow}$
      ]
    ]
    [.$\bnfnt{verb-phrase}$
      [.$\bnfnt{verb}$
        $\bnft{jumped}$
      ]
      [.$\bnfnt{prep-phrase}$
        [.$\bnfnt{prep}$
          $\bnft{over}$
        ]
        [.$\bnfnt{noun-phrase}$
          [.$\bnfnt{article}$
            $\bnft{the}$
          ]
          [.$\bnfnt{noun}$
            $\bnft{moon}$
          ]
        ]
      ]
    ]
  ]

\section{Backus-Naur Form}

A formal grammar is defined over a fixed collection of symbols.
This collection is divided into two disjoint groups: the \textbf{terminal symbols} and the \textbf{non-terminal} symbols.
We denote nonterminal symbols using notation of the form $\bnfnt{non-terminal}$, where \enquote{\textsf{non-terminal}} is replaced by an identifier.
We denote terminal symbols in red typewriter font like $\bnft{terminal}$.
%
\begin{remark}
We almost never state outright what our underlying symbols are.
We can determine the terminal and non-terminal symbols of a grammar by looking at its specification, assuming that every symbol must appear at least once in the specification.
\end{remark}
%
In the derivation from the previous section, we built a sequence of not-quite sentences, until the very last one, which was the sentence we wanted to prove grammatically correct.
We call these possibly-not-quite sentences \textbf{sentential forms}.
%
\begin{definition}
  A \textbf{sentential form} is a sequence of symbols (terminal or non-terminal).
  A \textbf{sentence} is a sequence of terminal symbols.
  We will use the symbol $\bnfempty$ to refer to the empty sentential form.
\end{definition}
%
\begin{remark}
We denote sentential forms using whitespace as the delimiter.
But it's important to recognize that \textit{this is just notation.}
We're not interested in low-level concerns like whitespace when studying formal grammar.
It may be useful to think of a sentential form as a \textit{list} of symbols.
\end{remark}
%
\begin{example}
The following are sentential forms that appear above:
\begin{align*}
  & \bnfnt{sentence} \\
  & \bnfnt{noun-phrase} \ \bnfnt{verb} \ \bnfnt{prep} \ \bnfnt{article} \ \bnfnt{noun} \\
  & \bnft{the} \ \bnft{cow} \ \bnft{jumped} \ \bnft{over} \ \bnft{the} \ \bnft{moon}
\end{align*}
but only the last one is a sentence.
\end{example}
In the process of constructing sentential forms, we replaced a non-terminal symbol with a sentential form, e.g., we replaced $\bnfnt{noun-phrase}$ with $\bnfnt{article} \ \bnfnt{noun}$.
These replacements are specified using \textbf{production rules}.
%
\begin{definition}
A \textbf{production rule} is a nonterminal symbol together with a sentential form.
Production rules are denoted using equations of the form
\begin{align*}
\bnfnt{non-terminal} \bnfeq \textsf{SENTENTIAL-FORM}
\end{align*}
where the left-hand side is the nonterminal symbol and the right-hand side is the sentential form.
\end{definition}
We read a production rule as: \enquote{the non-terminal symbol on the left-hand side can be replaced with the sentential form on the right hand side.}
%
\begin{example}
We can express that a prepositional phrase can be made up of a preposition followed by a noun-phrase using the following production rule:
%
\begin{align*}
\bnfnt{prep-phrase} \bnfeq \bnfnt{prep} \ \bnfnt{noun-phrase}
\end{align*}
\end{example}
%
\textbf{Backus-Naur Form (BNF) specifications} are used to describe \textbf{context-free grammars}, which form a class of formal grammars that are expressive enough to capture the syntax of most programming languages.
A BNF grammar is determined by what replacements we're allowed to perform and what kind of things we're trying to derive.
%
\begin{definition}
  A \textbf{Backus-Naur Form (BNF) specification} is a collection of production rules, together with a designated symbol called the \textbf{starting symbol}.
  If the start symbols is not explicitely given, it is taken to be the left-hand side of the \textit{first} rule appearing in the specification.
\end{definition}
%
\begin{remark}
  John Backus and Peter Naur are computer scientist who formalized this metasyntax in the process of developing and implementing the programming language \href{https://dl.acm.org/doi/pdf/10.1145/366193.366201}{ALGOL 60}.
\end{remark}
%
\begin{example}
  \label{simple-grammar}
  The following is a BNF specification for a grammar from the previous section.
  \begin{align*}
    \bnfnt{sentence} &\bnfeq \bnfnt{noun-phrase} \ \bnfnt{verb-phrase} \\
    \bnfnt{verb-phrase} &\bnfeq \bnfnt{verb} \ \bnfnt{prep-phrase} \\
    \bnfnt{verb-phrase} &\bnfeq \bnfnt{verb} \\
    \bnfnt{prep-phrase} &\bnfeq \bnfnt{prep} \ \bnfnt{noun-phrase} \\
    \bnfnt{noun-phrase} &\bnfeq \bnfnt{article} \ \bnfnt{noun} \\
    \bnfnt{article} &\bnfeq \bnft{the} \\
    \bnfnt{noun} &\bnfeq \bnft{cow} \\
    \bnfnt{noun} &\bnfeq \bnft{moon} \\
    \bnfnt{verb} &\bnfeq \bnft{jumped} \\
    \bnfnt{prep} &\bnfeq \bnft{over}
  \end{align*}
  The nonterminal symbols of this grammar are: $\bnfnt{sentence}$, $\bnfnt{noun-phrase}$, $\bnfnt{verb-phrase}$, $\bnfnt{verb}$, $\bnfnt{prep}$, $\bnfnt{prep-phrase}$, $\bnfnt{article}$, and $\bnfnt{noun}$.
  The terminal symbols of this grammar are: $\bnft{the}$, $\bnft{cow}$, $\bnft{moon}, $\bnft{jumped}, and $\bnft{over}$.
  The nonterminal symbol $\bnfnt{sentence}$ is our starting symbol (because it appears as the left-hand side of the first rule).
  This is, of course, not a complete or accurate model of English grammar (such a model would be incredibly complicated).
\end{example}
%
The last piece of our thought experiment that we need to formalize is the \textit{proof} of grammatical correction.
%
\begin{definition}
A \textbf{derivation} of a sentence $S$ in a BNF grammar $\mathcal G$ is a sequence of sentential forms $S_1, \dots, S_k$ with the following properties:
\begin{itemize}
  \item $S_1$ is the starting symbol of $\mathcal G$;
  \item $S_k = S$;
  \item for every index $i$ satisfying $i > 1$, the sentential form $S_i$ is the result of replacing a single nonterminal symbol in $S_{i - 1}$ with a sentential form, according to a production rule of $\mathcal G$.
\end{itemize}
We say that a grammar \textbf{recognizes} a sentence $S$ if there is a derivation of $S$ in the grammar.
\end{definition}
%
\begin{example}
The derivation from the previous section is a derivation is the sense of the previous definition.
We can also derive the sentence $\bnft{the} \ \bnft{moon} \ \bnft{jumped}$ as follows:
%
\begin{align*}
& \bnfnt{sentence} \\
& \bnfnt{noun-phrase} \ \bnfnt{verb-phrase} \\
& \bnfnt{article} \ \bnfnt{noun} \ \bnfnt{verb-phrase} \\
& \bnft{the} \ \bnfnt{noun} \ \bnfnt{verb-phrase} \\
& \bnft{the} \ \bnft{moon} \ \bnfnt{verb-phrase} \\
& \bnft{the} \ \bnft{moon} \ \bnfnt{verb} \\
& \bnft{the} \ \bnft{moon} \ \bnft{jumped}
\end{align*}
%
\end{example}
%
\begin{exercise}
For the derivation in the previous example, determine which production
rule is applied in each line of the derivation.
\end{exercise}
%
A sentence is not guaranteed to have a \textit{unique} derivation, but there a form of derivation we will single out.
%
\begin{definition}
A \textbf{leftmost derivation} of a sentence is one in which the leftmost nonterminal symbol is expanded in each step.
\end{definition}
%
The derivation given in the previous example is a leftmost derivation,
whereas the derivation in the previous section is not.
%
\begin{exercise}
Write a leftmost derivation of the sentence $\bnft{the} \ \bnft{cow} \ \bnft{jumped} \ \bnft{over} \ \bnft{the} \ \bnft{moon}$ in the grammar from \Cref{simple-grammar}.
\end{exercise}
%
\begin{remark}
In addition to context-free grammars, there are also context-\textit{sensitive} grammars.
In a context-sensitive grammar, production rules can transform \textit{sequences of} symbols.
In other words, for context-sensitive grammars, the left-hand side of a production rule is not required to be a single symbol.
\end{remark}
%
As we've noted several times now, Grammars imbue sentences with hierarchical structure.
We can represent this structure formally in terms of \textbf{parse trees}.
%
\begin{definition}
A \textbf{parse tree} for a sentence $S$ recognized by a grammar $\mathcal G$ is a tree $T$ with the following properties:
\begin{itemize}
\item the root of $T$ is the starting symbol of $\mathcal G$;
\item For every node $n$ of $T$, if $n$ has children $T_1, \dots, T_k$ then
\begin{align*}
\mathsf{val}(n) \bnfeq \mathsf{root}(T_1) \ \dots \mathsf{root}(T_k)
\end{align*}
is a production rule of $\mathcal G$;
\item The leaves in order from left to right (i.e., the \textbf{frontier} of $T$) form the sentence $S$.
In particular, the leaves of $T$ are labeled with terminal symbols of $\mathcal G$.
\end{itemize}
\end{definition}
%
\begin{example}
The tree at the end of the previous section is a parse tree for the sentence $\bnft{the} \ \bnft{cow} \ \bnft{jumped} \ \bnft{over} \ \bnft{the} \ \bnft{moon}$ in the grammar from \Cref{simple-grammar}.  The following is a parse tree for the sentence $\bnft{the} \ \bnft{moon} \ \bnft{jumped}$ in the same grammar:

\Tree
  [.$\bnfnt{sentence}$
    [.$\bnfnt{noun-phrase}$
      [.$\bnfnt{article}$
        $\bnft{the}$
      ]
      [.$\bnfnt{noun}$
        $\bnft{moon}$
      ]
    ]
    [.$\bnfnt{verb-phrase}$
      [.$\bnfnt{verb}$
        $\bnft{jumped}$
      ]
    ]
  ]
\end{example}
%
If a sentence $S$ has a derivation in a grammar $\mathcal G$, then it also has a parse tree, and vice versa.
Multiple derivations may correspond to the same parse tree, but to every parse tree there is a unique corresponding leftmost derivation; a leftmost derivation corresponds to a pre-order traversal of a parse tree.
%
\begin{exercise}
Implement the OCaml function
%
\begin{ocaml}
val frontier : 'a tree -> 'a list
\end{ocaml}
%
where \texttt{frontier t} is the list the values at the leaves of \texttt{t} in order from left to right.
You should use the following ADT for trees:
%
\begin{ocaml}
type 'a tree =
  | Leaf of 'a
  | Node of 'a * 'a tree list
\end{ocaml}
%
\end{exercise}

\section{Examples}

In this short section we look at a couple of examples that are closer to real-world examples we'll see.


\section{Grammatical Ambiguity}

\section{Extended BNF}

\section{Regular Expressions}
